%% 2-5 比・割合・百分率（12ページ、p.62〜73）第7回：比を2-4から移した章
%% 第11回：ページ制限なし。比の値・等しい比・比のつなぎ方・割合の3つの式・百分率⇔割合に「なぜ」の段を足し、  % noruby
%%   例題（歯の数から回転をつなぐ・100%をこえる割合・こう配）を足した。囲みの小見出しは1行にした（X19）。  % noruby
%% この章の部品（名前は \bkiif…）
\makeatletter
% 歯車（第5巻 5-6 の \bkgear と同じ形）  \bkiifgear{(中心)}{歯数}{モジュール mm}{1つの歯の向き（度）}{ぬり}{線}
%   かみ合う2つは、中心どうしを モジュール×(歯数1＋歯数2)÷2 はなし、中心線の上で「片方は歯、もう片方は歯みぞ」にする
\newcommand\bkiifgear[6]{\begin{scope}[shift={#1}]
  \pgfmathsetmacro\bkiif@R{#3*#2/2}\pgfmathsetmacro\bkiif@A{\bkiif@R+#3}%
  \pgfmathsetmacro\bkiif@F{\bkiif@R-1.25*#3}\pgfmathsetmacro\bkiif@P{360/#2}%
  \pgfmathsetmacro\bkiif@T{27*#3/\bkiif@A}\pgfmathsetmacro\bkiif@B{54*#3/\bkiif@F}%
  \filldraw[fill=#5, draw=#6, line width=0.5pt, line join=round]
    (#4-\bkiif@B:\bkiif@F)
    \foreach \k [evaluate=\k as \c using #4+\k*\bkiif@P] in {0,...,\numexpr#2-1\relax}{
      -- (\c-\bkiif@B:\bkiif@F) -- (\c-\bkiif@T:\bkiif@A) -- (\c+\bkiif@T:\bkiif@A) -- (\c+\bkiif@B:\bkiif@F)
      arc[start angle=\c+\bkiif@B, end angle=\c+\bkiif@P-\bkiif@B, radius=\bkiif@F]} -- cycle;
  \filldraw[fill=white, draw=#6, line width=0.5pt] (0,0) circle ({max(0.9,0.22*\bkiif@R)});
  \fill[#6] (0,0) circle (0.55);
\end{scope}}
% ジュースのコップ  \bkiifcup{(底のまん中)}{原液の大さじの数}{水の大さじの数}{ジュースの色}
%   左に大さじ1ぱいを1マスにしたぼう（だいだい＝原液、水色＝水）。ジュースの高さは、ぼうの高さと同じ
\newcommand\bkiifcup[4]{\begin{scope}[shift={#1}]
  \foreach \k in {1,...,#2} {\filldraw[fill=bkOrange!75, draw=white, line width=0.6pt] (-17,{(\k-1)*3.2}) rectangle ++(6,3.2);}
  \foreach \k in {1,...,#3} {\filldraw[fill=bkBlue!40, draw=white, line width=0.6pt] (-17,{(#2+\k-1)*3.2}) rectangle ++(6,3.2);}
  \pgfmathsetmacro\bkiif@h{(#2+#3)*3.2}\pgfmathsetmacro\bkiif@w{6.4+1.6*\bkiif@h/36}
  \fill[#4] (-6.4,0) -- (6.4,0) -- (\bkiif@w,\bkiif@h) -- (-\bkiif@w,\bkiif@h) -- cycle;
  \draw[bkInk!75, line width=1pt, line join=round] (-8,36) -- (-6.4,0) -- (6.4,0) -- (8,36);
\end{scope}}
% おかしのふくろ  \bkiifbag{(下のまん中)}{はば}{高さ}{色}{文字}
\newcommand\bkiifbag[5]{\begin{scope}[shift={#1}]\colorlet{bkiif@bc}{#4}%
  \fill[bkiif@bc, rounded corners=1.2mm] (-#2/2,0) rectangle (#2/2,#3);
  \fill[bkiif@bc!70!black] (-#2/2,#3-2.4) rectangle (#2/2,#3);
  \foreach \i in {1,...,9} {\fill[bkiif@bc!70!black] ({-#2/2+(\i-0.5)*#2/9-0.6},#3) -- ({-#2/2+(\i-0.5)*#2/9},#3+1) -- ({-#2/2+(\i-0.5)*#2/9+0.6},#3) -- cycle;}
  \fill[white, rounded corners=1mm] (-#2/2+2.5,#3*0.32) rectangle (#2/2-2.5,#3*0.68);
  \node[font=\sffamily\bfseries\footnotesize, text=bkInk, inner sep=0pt] at (0,#3*0.5) {#5};
\end{scope}}
% 警戒標識（上りのこう配）  \bkiifsign{(中心)}{大きさ}
\newcommand\bkiifsign[2]{\begin{scope}[shift={#1}, scale=#2]
  \fill[bkInk!80] (0,-0.3) rectangle (0.9,-12);
  \filldraw[fill=bkYellow, draw=bkInk, line width=0.9pt, rounded corners=0.8mm] (0,9) -- (9,0) -- (0,-9) -- (-9,0) -- cycle;
  \draw[bkInk, line width=0.5pt, rounded corners=0.6mm] (0,7.8) -- (7.8,0) -- (0,-7.8) -- (-7.8,0) -- cycle;
  \fill[bkInk] (-4.2,-3.2) -- (4.2,-3.2) -- (4.2,1.6) -- cycle;
  \node[font=\sffamily\bfseries\fontsize{8}{8}\selectfont, text=bkInk, inner sep=0pt] at (-1.6,1.6) {10\%};
\end{scope}}
\makeatother

\bkchapter{2}{5}{\R{比}{ひ}・\R{割合}{わり|あい}・\R{百分率}{ひゃく|ぶん|りつ}}{\bklevelrange{1}{2}}{\R{比}{ひ}・\R{割合}{わり|あい}・\R{百分率}{ひゃく|ぶん|りつ}}
\bkmichishirube{1-14,2-3,2-4}{2-6,2-15,5-6}

\begin{goal}
\begin{itemize}[label={\color{bkTheme}●}, leftmargin=1.6\zw, itemsep=0pt, topsep=1pt]
\item \bklevel{1}\enspace 2つの\R{量}{りょう}の\R{関係}{かん|けい}を\R{比}{ひ}で\R{表}{あらわ}し、\R{等}{ひと}しい\R{比}{ひ}をつくれる。
\item \bklevel{2}\enspace \R{比}{ひ}をかけ\R{算}{ざん}でつないで、\R{歯車}{は|ぐるま}の\R{回転}{かい|てん}を\R{求}{もと}められる。
\item \bklevel{1}\enspace \R{割合}{わり|あい}を\R{百分率}{ひゃく|ぶん|りつ}で\R{表}{あらわ}し、「20\%\R{引}{び}き」の\R{量}{りょう}を\R{求}{もと}められる。
\end{itemize}
\end{goal}

\begin{hatena}
テストで\R{全部}{ぜん|ぶ}\R{正解}{せい|かい}すると、\R{正}{ただ}しく\R{答}{こた}えた\R{割合}{わり|あい}は「100\%」。それより\R{上}{うえ}はありません。

ところが、お\R{店}{みせ}で「いつもの20\%\R{増量}{ぞう|りょう}」というおかしを\R{見}{み}つけました。いつものふくろとくらべると、\R{何}{なん}\%になるのかな？　100\%をこえることは、あるのでしょうか。

\begin{bkfigure}
\begin{tikzpicture}[x=1mm, y=1mm]
  % テストの紙
  \fill[white, draw=bkGray, line width=0.8pt] (0,0) rectangle (30,36);
  \foreach \y in {26,20,14,8} {\draw[bkLine, line width=0.8pt] (4,\y) -- (22,\y);
    \draw[bkRed, line width=1.1pt, line cap=round] (23.5,\y+0.4) -- (25,\y-1.2) -- (27.6,\y+2.6);}
  \node[font=\sffamily\bfseries\footnotesize, text=bkInk, inner sep=0pt] at (12,31.5) {テスト};
  \node[font=\sffamily\bfseries\small, text=bkRed, inner sep=0pt] at (15,-5) {100\%};
  % いつものふくろと増量のふくろ
  \bkiifbag{(64,0)}{20}{26}{bkBlue!55}{いつも}
  \bkiifbag{(97,0)}{22}{31.2}{bkOrange!80}{\R{増量}{ぞう|りょう}}
  \fill[bkRed] ($(113,33)+(0:8.2)$) \foreach \k in {1,...,23} { -- ($(113,33)+(\k*15:{7.2+cos(\k*180)})$)} -- cycle;
  \node[font=\sffamily\bfseries\scriptsize, text=white, inner sep=0pt] at (113,33) {+20\%};
  \node[font=\sffamily\bfseries\small, text=bkBlue!80!black, inner sep=0pt] at (64,-5) {100\%};
  \node[font=\sffamily\bfseries\small, text=bkOrange!80!black, inner sep=0pt] at (97,-5) {\R{何}{なん}\%？};
\end{tikzpicture}
\end{bkfigure}
\end{hatena}

\clearpage
%% ---------------- p.63 やってみよう・思い出そう ----------------
\begin{yattemiyou}[\R{同}{おな}じ\R{味}{あじ}のジュースをつくろう]
\par\addvspace{3pt}\textbf{\R{用意}{よう|い}するもの}\par\noindent うすめて\R{飲}{の}むジュースの\R{原液}{げん|えき}、\R{水}{みず}、\R{大}{おお}さじ、コップ3つ
\begin{enumerate}[label={\bkcircled[bkGreen]{\arabic*}}, leftmargin=2.4\zw, itemsep=1pt, topsep=3pt]
\item コップAに、\R{原液}{げん|えき}を\R{大}{おお}さじ1ぱい、\R{水}{みず}を\R{大}{おお}さじ4はい\R{入}{い}れて、まぜよう。
\item コップBには\R{原液}{げん|えき}2はいと\R{水}{みず}8ぱい、Cには\R{原液}{げん|えき}2はいと\R{水}{みず}4はいを\R{入}{い}れよう。
\item \R{飲}{の}みくらべよう。\R{同}{おな}じ\R{味}{あじ}のコップは、どれとどれかな？　いちばんこいのは、どれかな？
\end{enumerate}
\begin{bkfigure}
\begin{tikzpicture}[x=1mm, y=1mm]
  \bkiifcup{(22,0)}{1}{4}{bkOrange!30}
  \bkiifcup{(64,0)}{2}{8}{bkOrange!30}
  \bkiifcup{(106,0)}{2}{4}{bkOrange!58}
  \foreach \x/\n/\a/\b in {22/A/1/4, 64/B/2/8, 106/C/2/4} {
    \node[circle, fill=bkGreen, text=white, font=\sffamily\bfseries\small, minimum size=5.4mm, inner sep=0pt] at (\x-11,-5.5) {\n};
    \node[anchor=west, font=\sffamily\footnotesize, text=bkInk, inner sep=0pt] at (\x-7,-5.5) {\R{原液}{げん|えき}\a\,・\,\R{水}{みず}\b};}
  \fill[bkOrange!75] (40,38.5) rectangle ++(4,3.2);
  \node[anchor=west, font=\sffamily\footnotesize, text=bkInk, inner sep=0pt] at (45.5,40.1) {\R{原液}{げん|えき}\R{大}{おお}さじ1ぱい};
  \fill[bkBlue!40] (81,38.5) rectangle ++(4,3.2);
  \node[anchor=west, font=\sffamily\footnotesize, text=bkInk, inner sep=0pt] at (86.5,40.1) {\R{水}{みず}\R{大}{おお}さじ1ぱい};
\end{tikzpicture}
\end{bkfigure}
\end{yattemiyou}

\begin{omoidasou}
\bkrecall{2-3}{1あたりの\R{量}{りょう}は、わり\R{算}{ざん}でつくる。\R{速}{はや}さ＝\R{道}{みち}のり÷\R{時間}{じ|かん}}
\bkrecall{2-4}{\R{分数}{ぶん|すう}はわり\R{算}{ざん}の\R{答}{こた}え：$1\div4=\dfrac{1}{4}=0.25$。\R{分母}{ぶん|ぼ}と\R{分子}{ぶん|し}を\R{同}{おな}じ\R{数}{かず}でわっても、\R{大}{おお}きさは\R{同}{おな}じ（\R{約分}{やく|ぶん}）}
\bkrecall{2-4}{0と1のあいだの\R{数}{かず}をかけると、もとより\R{小}{ちい}さくなる：$1000\times0.8=800$}
\end{omoidasou}

\bkfukidashi{mina}{AとBは\R{同}{おな}じ\R{味}{あじ}だったよ。Bは\R{量}{りょう}が2\R{倍}{ばい}なのに、\R{味}{あじ}は\R{変}{か}わらない。Cはこかった！}

%% ---------------- 数学の道具 ----------------
\bksection{doguu}{\R{数学}{すう|がく}の\R{道具}{どう|ぐ}}

\bksubsection{1}{\R{比}{ひ}と\R{比}{ひ}の\R{値}{あたい}}

\leavevmode\label{bk:2-5:hi}コップAの\R{原液}{げん|えき}と\R{水}{みず}の\R{量}{りょう}の\R{関係}{かん|けい}を $1:4$ と\R{書}{か}き、「1\R{対}{たい}4」と\R{読}{よ}みます。このように\R{表}{あらわ}したものを\textbf{\R{比}{ひ}}といいます。\bkidx{ひ}{\R{比}{ひ}}Bは $2:8$、Cは $2:4$ です。

\begin{bkfigure}
\begin{tikzpicture}[x=1mm, y=1mm]
  \tikzset{gen/.style={fill=bkOrange!75, draw=white, line width=0.6pt}, miz/.style={fill=bkBlue!40, draw=white, line width=0.6pt}}
  % A
  \node[anchor=east, font=\sffamily\bfseries\small, text=bkInk, inner sep=0pt] at (7,11.5) {A};
  \fill[gen] (10,9) rectangle ++(5,5);
  \foreach \k in {1,...,4} {\fill[miz] (10+\k*5,9) rectangle ++(5,5);}
  \node[anchor=west, font=\sffamily\small, text=bkInk, inner sep=0pt] at (37.5,11.5) {$1:4$};
  % B（Aが2つ分）
  \node[anchor=east, font=\sffamily\bfseries\small, text=bkInk, inner sep=0pt] at (7,0) {B};
  \foreach \g in {0,1} {
    \fill[gen] (10+\g*28,-2.5) rectangle ++(5,5);
    \foreach \k in {1,...,4} {\fill[miz] (10+\g*28+\k*5,-2.5) rectangle ++(5,5);}
    \draw[bkGreen!70!black, line width=0.9pt, rounded corners=1mm] (9.2+\g*28,-3.3) rectangle ++(26.6,6.6);}
  \node[anchor=west, font=\sffamily\small, text=bkInk, inner sep=0pt] at (66,0) {$2:8$};
  \node[anchor=west, font=\sffamily\footnotesize, text=bkGreen!70!black, inner sep=0pt] at (78,0) {（Aが2つ\R{分}{ぶん}）};
  % C
  \node[anchor=east, font=\sffamily\bfseries\small, text=bkInk, inner sep=0pt] at (71,11.5) {C};
  \foreach \k in {0,1} {\fill[gen] (74+\k*5,9) rectangle ++(5,5);}
  \foreach \k in {2,...,5} {\fill[miz] (74+\k*5,9) rectangle ++(5,5);}
  \node[anchor=west, font=\sffamily\small, text=bkInk, inner sep=0pt] at (106.5,11.5) {$2:4$};
\end{tikzpicture}
\bkfigcaption{1マスが\R{大}{おお}さじ1ぱい。BはAを2つ\R{合}{あ}わせたものなので、\R{量}{りょう}は2\R{倍}{ばい}でも\R{味}{あじ}は\R{同}{おな}じ。}
\end{bkfigure}

$1:4$ の\R{前}{まえ}の\R{数}{かず}を\R{後}{うし}ろの\R{数}{かず}でわった $1\div4=0.25$ を、\textbf{\R{比}{ひ}の\R{値}{あたい}}といいます。\bkidx{ひのあたい}{\R{比}{ひ}の\R{値}{あたい}}\R{比}{ひ}の\R{値}{あたい}は、「\R{前}{まえ}の\R{量}{りょう}が、\R{後}{うし}ろの\R{量}{りょう}の\R{何倍}{なん|ばい}か」を\R{表}{あらわ}しています。Aでは、\R{原液}{げん|えき}は\R{水}{みず}の0.25\R{倍}{ばい}です。\R{分数}{ぶん|すう}で\R{書}{か}くと $1\div4=\dfrac{1}{4}$（2-4）で、\R{原液}{げん|えき}は\R{水}{みず}の $\dfrac{1}{4}$ だとわかります。

Bは $2\div8=\dfrac{2}{8}=\dfrac{1}{4}$（\R{約分}{やく|ぶん}）で、Aと\R{同}{おな}じです。だから\R{同}{おな}じ\R{味}{あじ}になりました。Cは $2\div4=\dfrac{1}{2}=0.5$ で\R{大}{おお}きく、\R{原液}{げん|えき}の\R{多}{おお}いこいジュースです。

\R{比}{ひ}の\R{値}{あたい}が\R{同}{おな}じ\R{比}{ひ}を\textbf{\R{等}{ひと}しい\R{比}{ひ}}といい、$1:4=2:8$ と\R{書}{か}きます。\R{図}{ず}のように、BはAを2つ\R{合}{あ}わせたものなので、\R{原液}{げん|えき}も\R{水}{みず}も2\R{倍}{ばい}です。このように、\R{比}{ひ}の\R{両方}{りょう|ほう}に\R{同}{おな}じ\R{数}{かず}をかけても、\R{比}{ひ}の\R{値}{あたい}は\R{変}{か}わりません（2-4：\R{分子}{ぶん|し}と\R{分母}{ぶん|ぼ}に\R{同}{おな}じ\R{数}{かず}をかけても、\R{分数}{ぶん|すう}の\R{大}{おお}きさは\R{変}{か}わらない）。\R{同}{おな}じ\R{数}{かず}でわっても\R{同}{おな}じです。
\[ 1:4=2:8=3:12\qquad\Bigl(\frac{1}{4}=\frac{2}{8}=\frac{3}{12}\Bigr) \]
できるだけ\R{小}{ちい}さな\R{整数}{せい|すう}の\R{比}{ひ}にすることを、「\R{比}{ひ}をかんたんにする」といいます。$12:18$ なら、\R{両方}{りょう|ほう}を6でわって $2:3$ です。6は、12と18を\R{素数}{そ|すう}のかけ\R{算}{ざん}に\R{分}{わ}けると（2-4）、$12=2\times2\times3$、$18=2\times3\times3$ で、\R{両方}{りょう|ほう}にある $2\times3=6$ です。
\[ 12:18=(12\div6):(18\div6)=2:3 \]

\bksubsection{2}{\R{比}{ひ}はかけ\R{算}{ざん}でつなげる}

\leavevmode\label{bk:2-5:tsunagu}\R{図}{ず}の\R{歯車}{は|ぐるま}は、\R{歯}{は}の\R{数}{かず}がAは48、Bは16、Cは8です。かみ\R{合}{あ}ったところでは、2つの\R{歯車}{は|ぐるま}の\R{歯}{は}が\R{同}{おな}じ\R{数}{かず}ずつ\R{送}{おく}られます（2-4）。Aが1\R{回転}{かい|てん}すると\R{歯}{は}が48こ\R{送}{おく}られるので、Bは $48\div16=3$ \R{回転}{かい|てん}します。\R{回転}{かい|てん}の\R{比}{ひ}は $\mathrm{A}:\mathrm{B}=1:3$。\R{同}{おな}じように、Bが1\R{回転}{かい|てん}するとCは $16\div8=2$ \R{回転}{かい|てん}するので、$\mathrm{B}:\mathrm{C}=1:2$ です。Aが1\R{回転}{かい|てん}すると、Cは\R{何}{なん}\R{回転}{かい|てん}するかな？

\begin{bkfigure}
\begin{tikzpicture}[x=1mm, y=1mm]
  % 歯数 A 48・B 16・C 8（第5巻 5-6 と同じ歯車）。モジュール 0.7 mm、中心距離 22.4 mm・8.4 mm
  \bkiifgear{(0,0)}{48}{0.7}{0}{bkBlue!22}{bkBlue!75!black}
  \bkiifgear{(22.4,0)}{16}{0.7}{191.25}{bkOrange!28}{bkOrange!75!black}
  \bkiifgear{(30.8,0)}{8}{0.7}{180}{bkGreen!25}{bkGreen!75!black}
  % 回る向き（かみ合うとなりどうしは逆向き）
  \draw[-{Stealth[length=4pt]}, bkInk!75, line width=0.9pt] (130:11) arc[start angle=130, end angle=50, radius=11];
  \draw[-{Stealth[length=3.5pt]}, bkInk!75, line width=0.8pt] ([shift={(22.4,0)}]62:8.6) arc[start angle=62, end angle=114, radius=8.6];
  \draw[-{Stealth[length=3.5pt]}, bkInk!75, line width=0.8pt] ([shift={(30.8,0)}]60:5.4) arc[start angle=60, end angle=-60, radius=5.4];
  \node[font=\sffamily\bfseries\small, anchor=north east, inner sep=0pt] at (-13.8,-13.8) {A};
  \node[font=\sffamily\bfseries\small, anchor=north, inner sep=0pt] at (22.4,-7.8) {B};
  \node[font=\sffamily\bfseries\small, anchor=west, inner sep=0pt] at (37.8,0) {C};
  % 回転の数の箱
  \node[bk node, minimum width=17mm] (a) at (57,6) {A\\[1pt]1\R{回転}{かい|てん}};
  \node[bk node, draw=bkOrange, minimum width=17mm] (b) at (81,6) {B\\[1pt]3\R{回転}{かい|てん}};
  \node[bk node, draw=bkGreen, minimum width=17mm] (c) at (105,6) {C\\[1pt]6\R{回転}{かい|てん}};
  \draw[bk link] (a) -- node[above, bk link label] {$\times3$} (b);
  \draw[bk link] (b) -- node[above, bk link label] {$\times2$} (c);
  \draw[bk link, bkTheme, rounded corners=2mm] (a.south) -- ++(0,-6) -| node[pos=0.25, below, bk link label, text=bkTheme!80!black] {$\times6\ (=\times3\times2)$} (c.south);
\end{tikzpicture}
\bkfigcaption{Bが3\R{回転}{かい|てん}すると、Cは $3\times2=6$ \R{回転}{かい|てん}。$\times3$ と $\times2$ をつないで $\times6$。}
\end{bkfigure}

\begin{enumerate}[label={\bkcircled[bkTheme]{\arabic*}}, leftmargin=2.4\zw, itemsep=1pt, topsep=2pt]
\item Aが1\R{回転}{かい|てん}すると、Bは3\R{回転}{かい|てん}します（×3）。
\item Bが1\R{回転}{かい|てん}するごとに、Cは2\R{回転}{かい|てん}します。Bが3\R{回転}{かい|てん}すると、Cは「2\R{回転}{かい|てん}」が3\R{回分}{かい|ぶん}で、$2\times3=6$ \R{回転}{かい|てん}します（×2）。
\end{enumerate}
\R{確}{たし}かめてみましょう。CはAの\R{歯}{は}48こ\R{分}{ぶん}を（Bを\R{通}{とお}して）\R{受}{う}けとるので、$48\div8=6$ \R{回転}{かい|てん}。\R{同}{おな}じになりました。\R{式}{しき}で\R{書}{か}くと、
\[ \mathrm{A}:\mathrm{B}=1:3,\quad \mathrm{B}:\mathrm{C}=1:2\quad\text{なら}\quad \mathrm{A}:\mathrm{C}=1:(3\times2)=1:6 \]
\R{比}{ひ}は、\textbf{かけ\R{算}{ざん}でつなげる}のです。Aを1\R{秒}{びょう}に2\R{回転}{かい|てん}させれば、Cは1\R{秒}{びょう}に $2\times6=12$ \R{回転}{かい|てん}します。

\bksubsection{3}{\R{割合}{わり|あい}：もとにする\R{量}{りょう}を1とみる}

\leavevmode\label{bk:2-5:wariai}シュートを10\R{本}{ぽん}\R{打}{う}って6\R{本}{ぽん}\R{入}{はい}ったAさんと、20\R{本}{ぽん}\R{打}{う}って11\R{本}{ぽん}\R{入}{はい}ったBさん。どちらが\R{上手}{じょうず}かな？　\R{打}{う}った\R{本数}{ほん|すう}がちがうので、1\R{本}{ぽん}あたりに\R{直}{なお}してくらべます。
\[ \text{Aさん}\quad 6\div10=0.6\qquad\qquad \text{Bさん}\quad 11\div20=0.55 \]

\begin{bkfigure}
\begin{tikzpicture}[x=1mm, y=1mm]
  % 2本の数直線（上が本数、下が割合）。1を90 mm
  \foreach \yy/\mo/\ku/\wa/\col/\nm in {14/10/6/0.6/bkBlue/Aさん, -6/20/11/0.55/bkOrange/Bさん} {
    \begin{scope}[shift={(0,\yy)}]
    \node[anchor=east, font=\sffamily\bfseries\small, text=\col!80!black, inner sep=0pt] at (-15,-3.5) {\nm};
    \node[anchor=east, font=\sffamily\footnotesize, text=bkInk, inner sep=0pt] at (-2.5,0) {\R{本数}{ほん|すう}};
    \node[anchor=east, font=\sffamily\footnotesize, text=bkInk, inner sep=0pt] at (-2.5,-7) {\R{割合}{わり|あい}};
    \draw[\col, line width=2.6pt] (0,0) -- (\wa*90,0);
    \draw[bkInk, line width=0.7pt] (0,0) -- (96,0);
    \draw[bkInk, line width=0.7pt] (0,-7) -- (96,-7);
    \foreach \t/\u in {0/0, \wa*90/\ku, 90/\mo} {\draw[bkInk, line width=0.7pt] (\t,-1.2) -- (\t,1.2);
      \node[anchor=south, font=\sffamily\footnotesize, inner sep=0pt] at (\t,2.2) {\u};}
    \foreach \t/\u in {0/0, \wa*90/\wa, 90/1} {\draw[bkInk, line width=0.7pt] (\t,-8.2) -- (\t,-5.8);
      \node[anchor=north, font=\sffamily\footnotesize, inner sep=0pt] at (\t,-9.2) {\u};}
    \draw[bkGray, line width=0.5pt, dash pattern=on 1.5pt off 1pt] (\wa*90,-1.6) -- (\wa*90,-5.4) (90,-1.6) -- (90,-5.4);
    \end{scope}}
\end{tikzpicture}
\bkfigcaption{\R{打}{う}った\R{本数}{ほん|すう}（もとにする\R{量}{りょう}）を1とみると、Aさんの0.6のほうが\R{大}{おお}きい。}
\end{bkfigure}

この0.6や0.55のように、もとにする\R{量}{りょう}を1とみたとき、くらべる\R{量}{りょう}がいくつにあたるかを\R{表}{あらわ}した\R{数}{かず}を\textbf{\R{割合}{わり|あい}}といいます。\bkidx{わりあい}{\R{割合}{わり|あい}}\R{数直線}{すう|ちょく|せん}から「10の0.6\R{倍}{ばい}が6」とわかるので、3つの\R{式}{しき}が\R{作}{つく}れます。

\begin{bkcenter}
\begin{bktable}{colspec={Q[l,m,wd=84mm] Q[l,m,wd=34mm]}, rowsep=1.5pt, row{1}={bg=white, font=\normalfont}}
\R{割合}{わり|あい}＝くらべる\R{量}{りょう}÷もとにする\R{量}{りょう} & $6\div10=0.6$ \\
くらべる\R{量}{りょう}＝もとにする\R{量}{りょう}×\R{割合}{わり|あい} & $10\times0.6=6$ \\
もとにする\R{量}{りょう}＝くらべる\R{量}{りょう}÷\R{割合}{わり|あい} & $6\div0.6=10$ \\
\end{bktable}
\bktabcaption{Aさんでたしかめた3つの\R{式}{しき}。\R{覚}{おぼ}えなくても、\R{数直線}{すう|ちょく|せん}をかけば\R{作}{つく}れる。}
\end{bkcenter}

3つの\R{式}{しき}は、それぞれ\R{次}{つぎ}のように\R{考}{かんが}えると\R{作}{つく}れます。
\begin{enumerate}[label={\bkcircled[bkTheme]{\arabic*}}, leftmargin=2.4\zw, itemsep=1pt, topsep=2pt]
\item \R{割合}{わり|あい}は「もとにする\R{量}{りょう}1あたり、くらべる\R{量}{りょう}はいくつか」。1あたりの\R{量}{りょう}なので、2-3の\R{速}{はや}さと\R{同}{おな}じくわり\R{算}{ざん}です（$6\div10=0.6$）。
\item \R{割合}{わり|あい}は「\R{何倍}{なん|ばい}か」でもあります。もとにする\R{量}{りょう}の10を0.6\R{倍}{ばい}すると、くらべる\R{量}{りょう}の6になります（$10\times0.6=6$）。
\item ②のかけ\R{算}{ざん}の\R{反対}{はん|たい}です。「\R{答}{こた}え×わる\R{数}{かず}＝わられる\R{数}{かず}」（2-4）なので、$6\div0.6=10$ でもとにする\R{量}{りょう}にもどります。
\end{enumerate}

\bksubsection{4}{\R{百分率}{ひゃく|ぶん|りつ}：もとを100とみる}

\noindent\begin{minipage}[c]{0.48\linewidth}
\leavevmode\label{bk:2-5:hyaku}\R{割合}{わり|あい}0.6は、もとを100とみると60にあたります。これを60\%と\R{書}{か}き、「60パーセント」と\R{読}{よ}みます。この\R{表}{あらわ}し\R{方}{かた}を\textbf{\R{百分率}{ひゃく|ぶん|りつ}}といいます。\bkidx{ひゃくぶんりつ}{\R{百分率}{ひゃく|ぶん|りつ}}
\end{minipage}\hfill
\begin{minipage}[c]{0.49\linewidth}\raggedleft
\bkkigocard{\%}{パーセント}{\R{百分率}{ひゃく|ぶん|りつ}の\R{記号}{き|ごう}。\R{割合}{わり|あい}の0.01が1\%}{percent}{2-5}
\end{minipage}

\R{野球}{や|きゅう}の\R{打率}{だ|りつ}「3\R{割}{わり}2\R{分}{ぶ}5\R{厘}{りん}」のように、0.1を1\R{割}{わり}、0.01を1\R{分}{ぶ}、0.001を1\R{厘}{りん}とする\R{表}{あらわ}し\R{方}{かた}を\textbf{\R{歩合}{ぶ|あい}}といいます。

\begin{bkcenter}
\begin{bktable}{colspec={Q[l,m,wd=20mm] Q[c,m,wd=15mm] Q[c,m,wd=15mm] Q[c,m,wd=27mm] Q[c,m,wd=14mm] Q[c,m,wd=14mm]}, rowsep=1.5pt}
\R{割合}{わり|あい} & 1 & 0.6 & 0.325 & 0.1 & 0.01 \\
\R{百分率}{ひゃく|ぶん|りつ} & 100\% & 60\% & 32.5\% & 10\% & 1\% \\
\R{歩合}{ぶ|あい} & 10\R{割}{わり} & 6\R{割}{わり} & {3\R{割}{わり}2\R{分}{ぶ}5\R{厘}{りん}} & 1\R{割}{わり} & 1\R{分}{ぶ} \\
\end{bktable}
\bktabcaption{\R{割合}{わり|あい}に100をかけると\R{百分率}{ひゃく|ぶん|りつ}}
\end{bkcenter}

なぜ100をかけるのでしょう。もとにする\R{量}{りょう}を「1」とみていたのを「100」とみなおすので、もとが100\R{倍}{ばい}になり、くらべる\R{量}{りょう}も100\R{倍}{ばい}になるからです。\R{反対}{はん|たい}に、\R{百分率}{ひゃく|ぶん|りつ}を\R{割合}{わり|あい}にもどすときは、100でわります（$30\%\to30\div100=0.3$、$115\%\to115\div100=1.15$）。\R{計算}{けい|さん}するときは、\R{割合}{わり|あい}（0.3や1.15）にもどしてから\R{使}{つか}います。

\bksubsection{5}{20\%\R{引}{び}きは0.8\R{倍}{ばい}、20\%\R{増}{ま}しは1.2\R{倍}{ばい}}

\leavevmode\label{bk:2-5:biki}1000\R{円}{えん}のくつが20\%\R{引}{び}きです。いくらかな？　もとのねだんを100\%とすると、はらうのは $100-20=80$（\%）で0.8\R{倍}{ばい}。$1000\times0.8=800$（\R{円}{えん}）です。ねびき $1000\times0.2=200$（\R{円}{えん}）を\R{引}{ひ}いても、\R{同}{おな}じになります。

\begin{bkfigure}
\begin{tikzpicture}[x=1mm, y=1mm]
  \foreach \y/\lab in {14/もと, 6/20\%\R{引}{び}き, -2/20\%\R{増量}{ぞう|りょう}} {\node[anchor=east, font=\sffamily\footnotesize, text=bkInk, inner sep=0pt] at (-2.5,\y+2.6) {\lab};}
  \foreach \k in {0,...,9} {\filldraw[fill=bkBlue!35, draw=white, line width=0.8pt] (\k*8,14) rectangle ++(8,5.2);}
  \foreach \k in {0,...,7} {\filldraw[fill=bkOrange!50, draw=white, line width=0.8pt] (\k*8,6) rectangle ++(8,5.2);}
  \foreach \k in {8,9} {\filldraw[fill=white, draw=bkGray, line width=0.6pt, dash pattern=on 1.2pt off 0.8pt] (\k*8+0.4,6.4) rectangle ++(7.2,4.4);}
  \foreach \k in {0,...,11} {\filldraw[fill=bkOrange!50, draw=white, line width=0.8pt] (\k*8,-2) rectangle ++(8,5.2);}
  \node[font=\sffamily\bfseries\footnotesize, inner sep=0pt] at (40,16.6) {1000\R{円}{えん}（100\%）};
  \node[font=\sffamily\bfseries\footnotesize, inner sep=0pt] at (32,8.6) {80\%};
  \node[font=\sffamily\footnotesize, text=bkGray!80!black, inner sep=0pt] at (72,8.6) {ねびき};
  \node[font=\sffamily\bfseries\footnotesize, inner sep=0pt] at (48,0.6) {120\%};
  \draw[bkInk!70, line width=0.6pt, dash pattern=on 1.5pt off 1pt] (80,-4) -- (80,21.5);
  \node[anchor=west, font=\sffamily\bfseries\footnotesize, text=bkOrange!80!black, inner sep=0pt] at (82,8.6) {0.8\R{倍}{ばい}};
  \node[anchor=west, font=\sffamily\bfseries\footnotesize, text=bkOrange!80!black, inner sep=0pt] at (98,0.6) {1.2\R{倍}{ばい}};
\end{tikzpicture}
\bkfigcaption{1マスが10\%。20\%\R{引}{び}きは8マス、20\%\R{増量}{ぞう|りょう}は12マス。}
\end{bkfigure}

20\%\R{増量}{ぞう|りょう}も、\R{同}{おな}じように\R{考}{かんが}えます。もとの100\%に20\%をたすと120\%で、\R{割合}{わり|あい}にもどすと1.2です。だから1.2\R{倍}{ばい}です。
\[ 100+20=120,\qquad 120\div100=1.2 \]

\R{第}{だい}1\R{巻}{かん}の\R{最後}{さい|ご}のカード「『20\%びき』ってどれだけ？」の\R{答}{こた}えは、「もとのねだんの0.8\R{倍}{ばい}」です。

\bksubsection{6}{100\%をこえる\R{割合}{わり|あい}、こえない\R{割合}{わり|あい}}

\leavevmode\label{bk:2-5:koeru}\bkhatenaref の「20\%\R{増量}{ぞう|りょう}」は、いつもを100\%とすると $100+20=120$（\%）です。くらべる\R{量}{りょう}がもとより\R{多}{おお}いので、\textbf{100\%をこえます}。「\R{去年}{きょ|ねん}の120\%」も\R{同}{おな}じです。

テストの\R{正解}{せい|かい}の\R{数}{かず}のように、くらべる\R{量}{りょう}がもとの\R{一部}{いち|ぶ}なら、100\%をこえません。モーターに\R{入}{い}れた\R{電気}{でん|き}のエネルギーも、\R{一部}{いち|ぶ}は\R{熱}{ねつ}や\R{音}{おと}になってにげます。\R{役}{やく}に\R{立}{た}つ\R{仕事}{し|ごと}になった\R{割合}{わり|あい}を\textbf{\R{効率}{こう|りつ}}といいます。エネルギーは、どこからかわいてきてふえることはないので、\R{効率}{こう|りつ}は100\%をこえません。\R{電気}{でん|き}のエネルギーは、1\R{秒}{びょう}あたりの\R{大}{おお}きさを W（ワット）という\R{単位}{たん|い}で\R{表}{あらわ}します。

\begin{motto}[2]
\textbf{3つの\R{比}{ひ}を1\R{度}{ど}に\R{書}{か}く}\par
\R{歯車}{は|ぐるま}の\R{回転}{かい|てん}は $\mathrm{A}:\mathrm{B}:\mathrm{C}=1:3:6$ とも\R{書}{か}けます。$\mathrm{A}:\mathrm{B}=2:3$、$\mathrm{B}:\mathrm{C}=4:5$ のように、Bの\R{数}{かず}がそろっていないときは、2-4の\R{通分}{つう|ぶん}と\R{同}{おな}じく、3と4の\R{公倍数}{こう|ばい|すう}の12にBをそろえます。$2:3$ の\R{両方}{りょう|ほう}に4、$4:5$ の\R{両方}{りょう|ほう}に3をかけます。
\[ \mathrm{A}:\mathrm{B}=8:12,\quad \mathrm{B}:\mathrm{C}=12:15\quad\text{だから}\quad \mathrm{A}:\mathrm{B}:\mathrm{C}=8:12:15 \]
\end{motto}

\begin{taisetsu}
\R{比}{ひ}の\R{両方}{りょう|ほう}に\R{同}{おな}じ\R{数}{かず}をかけても、わっても、\R{等}{ひと}しい\R{比}{ひ}になる。\\
\R{比}{ひ}はかけ\R{算}{ざん}でつなげる。$\mathrm{A}:\mathrm{B}=1:3$、$\mathrm{B}:\mathrm{C}=1:2$ なら $\mathrm{A}:\mathrm{C}=1:6$\\
\R{割合}{わり|あい}＝くらべる\R{量}{りょう}÷もとにする\R{量}{りょう}。0.01が1\%\\
20\%\R{引}{び}きは0.8\R{倍}{ばい}、20\%\R{増}{ま}しは1.2\R{倍}{ばい}
\end{taisetsu}

\begin{kotoba}
\small
\bkword{\R{比}{ひ}・\R{比}{ひ}の\R{値}{あたい}}{$1:4$（1\R{対}{たい}4）のように\R{表}{あらわ}した2つの\R{量}{りょう}の\R{関係}{かん|けい}。$1\div4$ が\R{比}{ひ}の\R{値}{あたい}。}
\bkword{\R{割合}{わり|あい}}{もとにする\R{量}{りょう}を1とみたとき、くらべる\R{量}{りょう}がいくつにあたるか。}
\bkword{\R{百分率}{ひゃく|ぶん|りつ}・\R{歩合}{ぶ|あい}}{もとを100とみた\R{割合}{わり|あい}（\%）。0.1を1\R{割}{わり}とする\R{表}{あらわ}し\R{方}{かた}。}
\bkword{\R{効率}{こう|りつ}}{\R{入}{い}れたエネルギーのうち、\R{役}{やく}に\R{立}{た}つ\R{仕事}{し|ごと}になった\R{割合}{わり|あい}。}
\end{kotoba}

\begin{bkfigure}
\begin{tikzpicture}[x=1mm, y=1mm]
  % 割合の3つの式をまとめた2本の数直線  % noruby
  \draw[bkInk, line width=0.8pt] (0,12) -- (100,12);
  \draw[bkInk, line width=0.8pt] (0,0) -- (100,0);
  \node[anchor=east, font=\sffamily\footnotesize] at (-2,12) {\R{量}{りょう}};
  \node[anchor=east, font=\sffamily\footnotesize] at (-2,0) {\R{割合}{わり|あい}};
  \foreach \x in {0,45,90} {\draw[bkInk, line width=0.7pt] (\x,10.5) -- (\x,13.5) (\x,-1.5) -- (\x,1.5);}
  \node[anchor=south, font=\small] at (0,14) {0};
  \node[anchor=south, font=\sffamily\footnotesize, text=bkBlue!80!black] at (45,14) {くらべる\R{量}{りょう}};
  \node[anchor=south, font=\sffamily\footnotesize, text=bkOrange!80!black] at (90,14) {もとにする\R{量}{りょう}};
  \node[anchor=north, font=\small] at (0,-2) {0};
  \node[anchor=north, font=\sffamily\footnotesize, text=bkBlue!80!black] at (45,-2) {\R{割合}{わり|あい}};
  \node[anchor=north, font=\small, text=bkOrange!80!black] at (90,-2) {1};
\end{tikzpicture}
\bkfigcaption{\R{上}{うえ}と\R{下}{した}が\R{同}{おな}じ\R{位置}{い|ち}どうしで\R{対応}{たい|おう}する。もとにする\R{量}{りょう}が1にあたる。}
\end{bkfigure}

\bkfukidashi{mina}{\R{割合}{わり|あい}の\R{問題}{もん|だい}は、まずこの2\R{本}{ほん}の\R{数直線}{すう|ちょく|せん}をかいて、どれがもとにする\R{量}{りょう}かを\R{決}{き}めるといいね。}

%% ---------------- 例題 ----------------
\bksection{reidai}{\R{例題}{れい|だい}でたしかめよう}

\begin{reidai}[\R{等}{ひと}しい\R{比}{ひ}]
（1）$24:36$ をかんたんにしましょう。（2）\R{酢}{す}と\R{油}{あぶら}を $2:3$ でまぜます。\R{酢}{す}が\qty{40}{mL}なら、\R{油}{あぶら}は\R{何}{なん}mL？\par
\kangaekata（1）は\R{両方}{りょう|ほう}を12でわる。（2）は\R{両方}{りょう|ほう}に $40\div2=20$ をかける。\par
\shiki（1）$24:36=(24\div12):(36\div12)=2:3$\par
\hspace*{2.6\zw}（2）$2:3=(2\times20):(3\times20)=40:60$\par
\kotae（1）$2:3$\quad（2）\qty{60}{mL}
\end{reidai}

\begin{jibun}
（1）$18:30$ をかんたんに。（2）\R{青}{あお}：\R{白}{しろ}＝$1:4$。\R{青}{あお}\qty{6}{g}なら\R{白}{しろ}は？
\bkjans{（1）\nobreak $3:5$\quad（2）\nobreak \qty{24}{g}}
\end{jibun}

\begin{reidai}[\R{比}{ひ}をかけ\R{算}{ざん}でつなぐ]
\R{歯車}{は|ぐるま}A・B・Cの\R{回転}{かい|てん}の\R{比}{ひ}は、$\mathrm{A}:\mathrm{B}=1:2$、$\mathrm{B}:\mathrm{C}=1:5$ です。Aを1\R{分}{ぷん}に10\R{回転}{かい|てん}させると、Cは1\R{分}{ぷん}に\R{何}{なん}\R{回転}{かい|てん}しますか。\par
\kangaekata $\times2$ と $\times5$ をつなぐと、$\mathrm{A}:\mathrm{C}=1:(2\times5)=1:10$\par
\shiki $10\times10=100$\kotaenext\mbox{100\R{回転}{かい|てん}}
\end{reidai}

\begin{jibun}
$\mathrm{A}:\mathrm{B}=1:4$、$\mathrm{B}:\mathrm{C}=1:3$ の\R{歯車}{は|ぐるま}。Aが1\R{秒}{びょう}に3\R{回転}{かい|てん}なら、Cは？
\bkjans{$\mathrm{A}:\mathrm{C}=1:12$ より、36\R{回転}{かい|てん}}
\end{jibun}

\begin{reidai}[\R{歯}{は}の\R{数}{かず}から\R{回転}{かい|てん}を\R{求}{もと}めてつなぐ]
\R{自転車}{じ|てん|しゃ}の\R{前}{まえ}のギアは\R{歯}{は}の\R{数}{かず}が36、\R{後}{うし}ろのギアは12です。\R{後}{うし}ろのギアとタイヤは\R{同}{おな}じ\R{軸}{じく}についていて、いっしょに\R{回}{まわ}ります。タイヤは1\R{回転}{かい|てん}で\qty{2}{m}\R{進}{すす}みます。ペダル（\R{前}{まえ}のギア）を1\R{分}{ぷん}に50\R{回転}{かい|てん}させると、1\R{分}{ぷん}で\R{何}{なん}m\R{進}{すす}みますか。\par
\kangaekata チェーンは、\R{前}{まえ}と\R{後}{うし}ろのギアの\R{歯}{は}を\R{同}{おな}じ\R{数}{かず}ずつ\R{送}{おく}る。ペダル→タイヤ（\R{歯}{は}の\R{数}{かず}でわる）と、タイヤ→\R{道}{みち}のり（\qty{2}{m}をかける）を、かけ\R{算}{ざん}でつなぐ。\par
\shiki ペダル1\R{回転}{かい|てん}で\R{歯}{は}が36こ\R{送}{おく}られるので、タイヤは $36\div12=3$（\R{回転}{かい|てん}）\par
\hspace*{2.6\zw}ペダル1\R{回転}{かい|てん}で\R{進}{すす}む\R{道}{みち}のり：$3\times2=6$（m）\par
\hspace*{2.6\zw}1\R{分}{ぷん}では：$50\times6=300$（m）\kotaenext\qty{300}{m}
\end{reidai}

\begin{jibun}
\R{前}{まえ}のギアの\R{歯}{は}が44、\R{後}{うし}ろのギアの\R{歯}{は}が22の\R{自転車}{じ|てん|しゃ}。タイヤは1\R{回転}{かい|てん}で\qty{2}{m}\R{進}{すす}みます。ペダルを1\R{分}{ぷん}に60\R{回転}{かい|てん}させると、1\R{分}{ぷん}で\R{何}{なん}m\R{進}{すす}みますか。
\bkjans{$44\div22=2$（\R{回転}{かい|てん}）、$2\times2=4$（m）、$60\times4=240$\quad\qty{240}{m}}
\end{jibun}

\begin{reidai}[\R{割合}{わり|あい}を\R{百分率}{ひゃく|ぶん|りつ}で]
（1）シュート25\R{本}{ほん}\R{中}{ちゅう}18\R{本}{ぽん}\R{入}{はい}った\R{割合}{わり|あい}は\R{何}{なん}\%？\R{歩合}{ぶ|あい}では？（2）\qty{50}{W}（ワット）の\R{電気}{でん|き}を\R{入}{い}れると、\qty{40}{W}\R{分}{ぶん}の\R{仕事}{し|ごと}をするモーターの\R{効率}{こう|りつ}は？\par
\kangaekata もとにする\R{量}{りょう}は、（1）は\R{打}{う}った25\R{本}{ほん}、（2）は\R{入}{い}れた\R{電気}{でん|き}です。\par
\shiki（1）$18\div25=0.72$\quad $0.72\times100=72$ で72\%\par
\hspace*{2.6\zw}\phantom{（1）}0.72は0.1が7こと0.01が2こなので、7\R{割}{わり}2\R{分}{ぶ}\par
\hspace*{2.6\zw}（2）$40\div50=0.8$\quad $0.8\times100=80$ で80\%\par
\kotae（1）72\%、7\R{割}{わり}2\R{分}{ぶ}\quad（2）80\%
\end{reidai}

\begin{jibun}
（1）40\R{問中}{もん|ちゅう}34\R{問}{もん}\R{正解}{せい|かい}は\R{何}{なん}\%？\R{歩合}{ぶ|あい}では？（2）\qty{200}{W}\R{入}{い}れて\qty{170}{W}\R{分}{ぶん}の\R{仕事}{し|ごと}。\R{効率}{こう|りつ}は？
\bkjans{（1）\nobreak 85\%、8\R{割}{わり}5\R{分}{ぶ}\quad（2）\nobreak $170\div200=0.85$ で85\%}
\end{jibun}

\begin{reidai}[100\%をこえる\R{割合}{わり|あい}]
\R{図書館}{と|しょ|かん}の\R{本}{ほん}を\R{借}{か}りた\R{人}{ひと}は、\R{去年}{きょ|ねん}が600\R{人}{にん}、\R{今年}{ことし}が690\R{人}{にん}でした。（1）\R{今年}{ことし}は\R{去年}{きょ|ねん}の\R{何}{なん}\%ですか。（2）\R{何}{なん}\%ふえましたか。\par
\kangaekata 「\R{去年}{きょ|ねん}の」とあるので、もとにする\R{量}{りょう}は\R{去年}{きょ|ねん}の600\R{人}{にん}。くらべる\R{量}{りょう}がもとより\R{多}{おお}いので、100\%をこえるはず。\par
\shiki（1）$690\div600=1.15$\quad $1.15\times100=115$（\%）\par
\hspace*{2.6\zw}（2）$115-100=15$（\%）\quad たしかめ：$600\times0.15=90$、$600+90=690$\par
\kotae（1）115\%\quad（2）15\%
\end{reidai}

\begin{jibun}
\R{去年}{きょ|ねん}250\R{本}{ぼん}とれたきゅうりが、\R{今年}{ことし}は300\R{本}{ぼん}とれました。（1）\R{今年}{ことし}は\R{去年}{きょ|ねん}の\R{何}{なん}\%ですか。（2）\R{何}{なん}\%ふえましたか。
\bkjans{（1）\nobreak $300\div250=1.2$ で120\%\quad（2）\nobreak 20\%}
\end{jibun}

\begin{reidai}[\R{百分率}{ひゃく|ぶん|りつ}から\R{量}{りょう}を\R{求}{もと}める]
（1）2500\R{円}{えん}のくつの30\%\R{引}{び}きは、いくら？（2）「20\%\R{増量}{ぞう|りょう}」で\qty{600}{mL}のジュース。いつもは\R{何}{なん}mL？\par
\kangaekata（1）はらうのは70\%。（2）もとにする\R{量}{りょう}＝くらべる\R{量}{りょう}÷\R{割合}{わり|あい}\par
\shiki（1）$2500\times0.7=1750$（\R{円}{えん}）\par
\hspace*{2.6\zw}（2）$600\div1.2=500$（mL）\quad たしかめ：$500\times1.2=600$\par
\kotae（1）1750\R{円}{えん}\quad（2）\qty{500}{mL}
\end{reidai}

\begin{jibun}
（1）1800\R{円}{えん}の15\%\R{引}{び}きは？（2）「25\%\R{増量}{ぞう|りょう}」で\qty{150}{g}。いつもは\R{何}{なん}g？（3）いつも\qty{450}{mL}のジュースが「20\%\R{増量}{ぞう|りょう}」になると\R{何}{なん}mL？
\bkjans{（1）\nobreak $1800\times0.85=1530$（\R{円}{えん}）\quad（2）\nobreak $150\div1.25=120$（g）\quad（3）\nobreak $450\times1.2=540$（mL）}
\end{jibun}

\begin{reidai}[\R{坂}{さか}のこう\R{配}{ばい}]
\R{水平}{すい|へい}に\qty{200}{m}\R{進}{すす}むと\qty{16}{m}\R{上}{あ}がる\R{坂}{さか}があります（こう\R{配}{ばい}は→\pageref{bk:2-5:koubai}ページ）。（1）こう\R{配}{ばい}は\R{何}{なん}\%ですか。（2）この\R{坂}{さか}を\R{水平}{すい|へい}に\qty{450}{m}\R{進}{すす}むと、\R{何}{なん}m\R{上}{あ}がりますか。\par
\kangaekata もとにする\R{量}{りょう}は\R{水平}{すい|へい}に\R{進}{すす}んだ\R{長}{なが}さ、くらべる\R{量}{りょう}は\R{上}{あ}がった\R{高}{たか}さ。（2）は、くらべる\R{量}{りょう}＝もとにする\R{量}{りょう}×\R{割合}{わり|あい}。\par
\shiki（1）$16\div200=0.08$\quad $0.08\times100=8$（\%）\par
\hspace*{2.6\zw}（2）$450\times0.08=36$（m）\par
\kotae（1）8\%\quad（2）\qty{36}{m}
\begin{bkfigure}
\begin{tikzpicture}[x=0.2mm, y=0.5mm]
  % こう配 8%：たては横の 2.5 倍にのばしてかいた  % noruby
  \fill[bkGreen!18] (0,0) -- (450,0) -- (450,36) -- cycle;
  \draw[bkInk, line width=1pt] (0,0) -- (450,36);
  \draw[bkInk!60, line width=0.6pt] (0,0) -- (450,0) -- (450,36);
  \draw[bkInk!50, line width=0.5pt, dash pattern=on 1.5pt off 1pt] (200,0) -- (200,16);
  \node[anchor=north, font=\sffamily\footnotesize] at (100,-2) {\qty{200}{m}};
  \node[anchor=north, font=\sffamily\footnotesize] at (325,-2) {\qty{250}{m}};
  \node[anchor=west, font=\sffamily\footnotesize] at (210,6) {\qty{16}{m}};
  \node[anchor=west, font=\sffamily\footnotesize] at (456,18) {\qty{36}{m}};
\end{tikzpicture}
\bkfigcaption{\R{水平}{すい|へい}に\qty{200}{m}で\qty{16}{m}、\qty{450}{m}で\qty{36}{m}。\R{高}{たか}さはたてにのばしてかいた。}
\end{bkfigure}
\end{reidai}

\begin{jibun}
\R{水平}{すい|へい}に\qty{300}{m}\R{進}{すす}むと\qty{15}{m}\R{上}{あ}がる\R{坂}{さか}。（1）こう\R{配}{ばい}は\R{何}{なん}\%？（2）\R{水平}{すい|へい}に\qty{800}{m}\R{進}{すす}むと、\R{何}{なん}m\R{上}{あ}がる？
\bkjans{（1）\nobreak $15\div300=0.05$ で5\%\quad（2）\nobreak $800\times0.05=40$\quad\qty{40}{m}}
\end{jibun}


%% ---------------- まちがえやすいところ ----------------
\bksection{tsumazuki}{まちがえやすいところをチェック}

\begin{tsumazuki}{20\%\R{引}{び}きを2\R{回}{かい}すると、40\%\R{引}{び}き？}
1000\R{円}{えん}を20\%\R{引}{び}きにして、さらにその20\%\R{引}{び}きにすると……
\bkng{$20+20=40$ で40\%\R{引}{び}き。$1000\times0.6=600$（\R{円}{えん}）}
\bkok{$1000\times0.8\times0.8=640$（\R{円}{えん}）。36\%\R{引}{び}き}
2\R{回}{かい}\R{目}{め}は800\R{円}{えん}がもとにする\R{量}{りょう}なので、\%どうしはたせません。\R{割合}{わり|あい}も、かけ\R{算}{ざん}でつなぎます。
\end{tsumazuki}

\begin{tsumazuki}{\R{逆}{ぎゃく}から\R{見}{み}ると、もとにする\R{量}{りょう}も\R{変}{か}わる}
\R{今年}{ことし}の180\R{人}{にん}は\R{去年}{きょ|ねん}の150\R{人}{にん}の120\%。では、\R{去年}{きょ|ねん}は\R{今年}{ことし}の\R{何}{なん}\%？
\bkng{$100-20=80$ で80\%}
\bkok{$150\div180=0.833\cdots$ で、およそ83\%（もとは\R{今年}{ことし}の180\R{人}{にん}）}
\end{tsumazuki}

\begin{tsumazuki}{\R{比}{ひ}の\R{順番}{じゅん|ばん}を\R{入}{い}れかえる}
「\R{原液}{げん|えき}1に\R{水}{みず}4」の\R{比}{ひ}は……
\bkng{$\text{\R{原液}{げん|えき}}:\text{\R{水}{みず}}=4:1$\quad（とてもこいジュースになる）}
\bkok{$\text{\R{原液}{げん|えき}}:\text{\R{水}{みず}}=1:4$}
\R{何}{なに}と\R{何}{なに}の\R{比}{ひ}かを、ことばで\R{書}{か}いてから\R{数}{かず}を\R{入}{い}れましょう。
\end{tsumazuki}

%% ---------------- 物理・工学ではこう使う ----------------
\begin{butsuri}
\begin{bkfigure}
\begin{tikzpicture}[x=1mm, y=1mm]
  % 自転車のギア：前は歯の数 42、後ろは歯の数 14（モジュール 0.5 mm）。チェーンはピッチ円の外側の接線
  \draw[bkInk!80, line width=1.6pt, line join=round]
    (1.598,10.378) -- (46.533,3.459) arc[start angle=81.25, end angle=-81.25, radius=3.5]
    -- (1.598,-10.378) arc[start angle=-81.25, end angle=-278.75, radius=10.5] -- cycle;
  \bkiifgear{(0,0)}{42}{0.5}{0}{bkGray!25}{bkInk!70}
  \bkiifgear{(46,0)}{14}{0.5}{0}{bkGray!25}{bkInk!70}
  \draw[bkInk!80, line width=1.6pt] (1.598,10.378) -- (46.533,3.459) (1.598,-10.378) -- (46.533,-3.459);
  % ペダル
  \draw[bkGray!80!black, line width=2pt, line cap=round] (0,0) -- (-5,-8);
  \fill[bkGray!70!black, rounded corners=0.5mm] (-8.5,-9.6) rectangle (-1.5,-6.8);
  \node[anchor=north, font=\sffamily\footnotesize, inner sep=0pt] at (2,-13.5) {\R{前}{まえ}　\R{歯}{は} 42};
  \node[anchor=north, font=\sffamily\footnotesize, inner sep=0pt] at (46,-6.5) {\R{後}{うし}ろ　\R{歯}{は} 14};
  % 表
  \node[anchor=west, inner sep=0pt] at (55,0) {\begin{tblr}{colspec={Q[c,m] Q[c,m] Q[c,m]}, rowsep=0.5pt, colsep=2.5pt,
     hlines={0.5pt,bkLine}, row{1}={font=\sffamily\bfseries\footnotesize, bg=bkPurple!12}, cells={font=\sffamily\footnotesize}}
     \R{後}{うし}ろの\R{歯}{は}の\R{数}{かず} & タイヤの\R{回転}{かい|てん} & \R{進}{すす}む\R{道}{みち}のり \\
     14 & $42\div14=3$ & \R{約}{やく}\qty{6}{m} \\
     21 & $42\div21=2$ & \R{約}{やく}\qty{4}{m} \\
     28 & $42\div28=1.5$ & \R{約}{やく}\qty{3}{m} \\
   \end{tblr}};
\end{tikzpicture}
\bkfigcaption{ペダル1\R{回転}{かい|てん}あたり。タイヤは1\R{回転}{かい|てん}で\R{約}{やく}\qty{2}{m}\R{進}{すす}む。}
\end{bkfigure}
\begin{itemize}[leftmargin=1.2\zw, itemsep=1pt, topsep=0pt, label={\color{bkPurple}●}]
\item \textbf{\R{自転車}{じ|てん|しゃ}の\R{変速}{へん|そく}ギア}\\
「ペダル→タイヤ」と「タイヤ→\R{道}{みち}」の\R{比}{ひ}を、かけ\R{算}{ざん}でつないでいます。のぼり\R{坂}{ざか}で\R{歯}{は}の\R{数}{かず}が28のギアにすると、\R{進}{すす}む\R{道}{みち}のりは\R{半分}{はん|ぶん}になるかわりに、ペダルが\R{軽}{かる}くなります。
\item \textbf{ロボットの\R{減速機}{げん|そく|き}}\\
モーターは\R{速}{はや}く\R{回}{まわ}るけれど、\R{回}{まわ}す\R{力}{ちから}は\R{弱}{よわ}いので、\R{歯車}{は|ぐるま}で\R{回転}{かい|てん}をおそくして\R{力}{ちから}を\R{強}{つよ}くします。1つの\R{軸}{じく}に\R{大小}{だい|しょう}2つの\R{歯車}{は|ぐるま}をつけて2\R{段}{だん}つなぐと、\R{回転}{かい|てん}は $\dfrac{1}{3}\times\dfrac{1}{4}=\dfrac{1}{12}$ のように、かけ\R{算}{ざん}で\R{大}{おお}きく\R{変}{か}わります。
\end{itemize}
\noindent\begin{minipage}[c]{0.64\linewidth}
\begin{itemize}[leftmargin=1.2\zw, itemsep=1pt, topsep=1pt, label={\color{bkPurple}●}]
\item \textbf{\R{坂}{さか}のこう\R{配}{ばい}}\\
\R{坂}{さか}の\R{急}{きゅう}さは、\R{水平}{すい|へい}に\R{進}{すす}んだ\R{長}{なが}さをもとにした、\R{上}{あ}がった\R{高}{たか}さの\R{割合}{わり|あい}で\R{表}{あらわ}します。これを\textbf{こう\R{配}{ばい}}といいます。\label{bk:2-5:koubai}\bkidx{こうばい}{こう\R{配}{ばい}}\R{水平}{すい|へい}に\qty{100}{m}で\qty{10}{m}\R{上}{あ}がる\R{坂}{さか}は、$10\div100=0.1$ で10\%。\R{道路}{どう|ろ}の\R{標識}{ひょう|しき}で\R{知}{し}らせます。\R{鉄道}{てつ|どう}では、もとを1000とみた‰（パーミル）を\R{使}{つか}い、\R{箱根}{はこ|ね}\R{登山}{と|ざん}\R{鉄道}{てつ|どう}には80‰の\R{坂}{さか}があります。
\end{itemize}
\end{minipage}\hfill
\begin{minipage}[c]{0.33\linewidth}\centering
\begin{tikzpicture}[x=1mm, y=1mm]
  \fill[bkGreen!20] (0,0) -- (30,0) -- (30,3) -- cycle;
  \draw[bkInk, line width=1pt] (0,0) -- (30,3);
  \draw[bkInk!60, line width=0.6pt] (0,0) -- (30,0) -- (30,3);
  \node[anchor=north, font=\sffamily\footnotesize, inner sep=0pt] at (15,-1.8) {\R{水平}{すい|へい}に\qty{100}{m}};
  \node[anchor=west, font=\sffamily\footnotesize, inner sep=0pt] at (31.5,1.5) {\qty{10}{m}};
  \bkiifsign{(12,18)}{0.85}
\end{tikzpicture}
\end{minipage}
\end{butsuri}

\begin{column}{コピー\R{機}{き}の「141\%」と「71\%」}
\noindent\begin{minipage}[c]{0.58\linewidth}
コピー\R{機}{き}で、A4の\R{紙}{かみ}をA3の\R{大}{おお}きさに\R{拡大}{かく|だい}するときは「141\%」、A3をA4に\R{縮小}{しゅく|しょう}するときは「71\%」をえらびます。もとにする\R{量}{りょう}は、もとの\R{紙}{かみ}の\R{辺}{へん}の\R{長}{なが}さです。A4のたての\R{長}{なが}さ\qty{297}{mm}の141\%は、$297\times1.41=418.77$ で\R{約}{やく}\qty{419}{mm}。A3のたての\R{長}{なが}さ\qty{420}{mm}とほぼ\R{同}{おな}じです。
\end{minipage}\hfill
\begin{minipage}[c]{0.38\linewidth}\centering
\begin{tikzpicture}[x=0.11mm, y=0.11mm]
  \filldraw[fill=bkBlue!10, draw=bkBlue!70!black, line width=0.8pt] (0,0) rectangle (297,420);
  \filldraw[fill=bkOrange!15, draw=bkOrange!80!black, line width=0.8pt] (0,0) rectangle (210,297);
  \node[font=\sffamily\bfseries\footnotesize] at (105,148) {A4};
  \node[font=\sffamily\bfseries\footnotesize] at (148,370) {A3};
\end{tikzpicture}
\end{minipage}
\par\vspace{3pt}
\R{反対}{はん|たい}に、A3の\R{長}{なが}さ\qty{420}{mm}の71\%は、$420\times0.71=298.2$ で\R{約}{やく}\qty{298}{mm}。A4の\qty{297}{mm}にほぼもどります。なぜ141と71という\R{半端}{はん|ぱ}な\R{数}{かず}なのかは、2-8の\R{平方根}{へい|ほう|こん}でわかります（1.41…×1.41…がちょうど2になる\R{数}{かず}）。
\end{column}

\begin{column}{スマートフォンの\R{電池}{でん|ち}の「\%」}
スマートフォンの\R{電池}{でん|ち}の\R{表示}{ひょう|じ}「25\%」は、\R{満}{まん}タンにためられる\R{電気}{でん|き}の\R{量}{りょう}をもとにする\R{量}{りょう}（100\%）にした\R{割合}{わり|あい}です。\R{満}{まん}タンで\R{動画}{どう|が}を10\R{時間}{じ|かん}\R{見}{み}られる\R{電池}{でん|ち}なら、25\%の\R{残}{のこ}りでは $10\times0.25=2.5$ で、\R{約}{やく}2.5\R{時間}{じ|かん}\R{見}{み}られます。\R{古}{ふる}くなった\R{電池}{でん|ち}は、\R{満}{まん}タンの\R{量}{りょう}そのもの（もとにする\R{量}{りょう}）が\R{小}{ちい}さくなるので、\R{同}{おな}じ「100\%」でも\R{短}{みじか}い\R{時間}{じ|かん}しか\R{使}{つか}えなくなります。
\begin{bkfigure}
\begin{tikzpicture}[x=1mm, y=1mm]
  \foreach \x/\p/\t in {0/1/{100\%}, 40/0.25/{25\%}} {
    \draw[bkInk!80, line width=0.9pt, rounded corners=1mm] (\x,0) rectangle (\x+24,10);
    \fill[bkInk!80] (\x+24,3) rectangle (\x+26,7);
    \fill[bkGreen!60, rounded corners=0.6mm] (\x+1.2,1.2) rectangle ({\x+1.2+21.6*\p},8.8);
    \node[anchor=west, font=\sffamily\bfseries\footnotesize] at (\x+28,5) {\t};}
\end{tikzpicture}
\end{bkfigure}
\end{column}

\bkfukidashi{sora}{チラシの「\R{全品}{ぜん|ぴん}30\%オフ」や\R{天気}{てん|き}\R{予報}{よ|ほう}の「\R{降水確率}{こう|すい|かく|りつ}」にも\%があるよ。\R{何}{なに}をもとにしているか、\R{考}{かんが}えてみよう！}





\clearpage
%% ---------------- 章末問題 ----------------
\bkendproblems
\bkproblemgroup{\R{基本}{き|ほん}}{1}

\bkproblem \R{比}{ひ}をかんたんにしましょう。\quad（1）\nobreak $8:12$\quad（2）\nobreak $45:30$\quad（3）\nobreak $36:48$
\bkpans{（1）\nobreak $2:3$\quad（2）\nobreak $3:2$\quad（3）\nobreak $3:4$}

\noindent\begin{minipage}[c]{0.64\linewidth}
\bkproblem \R{壁}{かべ}にぬるモルタルは、セメントと\R{砂}{すな}を $1:3$ の\R{比}{ひ}でまぜてつくります。セメントを\qty{4}{kg}\R{使}{つか}うと、\R{砂}{すな}は\R{何}{なん}kgいりますか。
\bkpans{\qty{12}{kg}}
\end{minipage}\hfill
\begin{minipage}[c]{0.33\linewidth}\centering
\begin{tikzpicture}[x=1mm, y=1mm]
  % セメント 1 ばいと砂 3 ばい（同じ大きさのバケツ）
  \foreach \k/\c in {0/bkGray!55, 1/bkYellow!55, 2/bkYellow!55, 3/bkYellow!55} {
    \fill[\c] ({\k*11},0) -- ({\k*11+8},0) -- ({\k*11+9},8) -- ({\k*11-1},8) -- cycle;
    \draw[bkInk!70, line width=0.6pt] ({\k*11-1},8) -- ({\k*11},0) -- ({\k*11+8},0) -- ({\k*11+9},8);}
  \node[anchor=north, font=\sffamily\footnotesize, inner sep=0pt] at (4,-1.8) {セメント};
  \draw[bkInk!60, line width=0.6pt, decorate, decoration={brace, amplitude=1.2mm}] (10,9.6) -- (42,9.6);
  \node[anchor=south, font=\sffamily\footnotesize, inner sep=0pt] at (26,12.4) {\R{砂}{すな}};
\end{tikzpicture}
\end{minipage}

\bkproblem \R{歯車}{は|ぐるま}A・B・Cの\R{回転}{かい|てん}の\R{比}{ひ}は、$\mathrm{A}:\mathrm{B}=1:3$、$\mathrm{B}:\mathrm{C}=1:4$ です。（1）$\mathrm{A}:\mathrm{C}$ を\R{求}{もと}めましょう。（2）Aを1\R{分}{ぷん}に20\R{回転}{かい|てん}させると、Cは1\R{分}{ぷん}に\R{何}{なん}\R{回転}{かい|てん}しますか。
\bkpans{（1）\nobreak $1:12$\quad（2）\nobreak 240\R{回転}{かい|てん}}

\bkproblem 20\R{本}{ぽん}のシュートのうち13\R{本}{ぼん}\R{入}{はい}りました。\R{入}{はい}った\R{割合}{わり|あい}を、\R{百分率}{ひゃく|ぶん|りつ}と\R{歩合}{ぶ|あい}で\R{表}{あらわ}しましょう。
\bkpans{65\%、6\R{割}{わり}5\R{分}{ぶ}}

\noindent\begin{minipage}[c]{0.64\linewidth}
\bkproblem （1）3200\R{円}{えん}のかばんの25\%\R{引}{び}きは、いくらですか。（2）「10\%\R{増量}{ぞう|りょう}」で\qty{440}{mL}のシャンプー。いつもは\R{何}{なん}mLですか。
\bkpans{（1）\nobreak 2400\R{円}{えん}\quad（2）\nobreak \qty{400}{mL}}
\end{minipage}\hfill
\begin{minipage}[c]{0.33\linewidth}\centering
\begin{tikzpicture}[x=1mm, y=1mm]
  % ねふだ
  \begin{scope}[shift={(0,8)}, rotate=-8]
    \filldraw[fill=white, draw=bkInk!70, line width=0.7pt] (0,0) -- (16,0) -- (16,12) -- (0,12) -- (-4,6) -- cycle;
    \draw[bkInk!70, line width=0.6pt] (-1.2,6) circle (0.9);
    \node[font=\sffamily\bfseries\footnotesize, inner sep=0pt] at (8,8.4) {3200\R{円}{えん}};
    \node[font=\sffamily\bfseries\footnotesize, text=bkRed, inner sep=0pt] at (8,3.2) {25\%\R{引}{び}き};
  \end{scope}
  % シャンプーのボトル
  \begin{scope}[shift={(30,0)}]
    \fill[bkGreen!25, rounded corners=1.5mm] (-7,0) rectangle (7,26);
    \draw[bkGreen!60!black, line width=0.8pt, rounded corners=1.5mm] (-7,0) rectangle (7,26);
    \fill[bkGray!60] (-1.4,26) rectangle (1.4,29.5);
    \fill[bkGray!70] (-1.4,29.5) rectangle (5.5,31);
    \fill[white] (-6,4) rectangle (6,11);
    \node[font=\sffamily\bfseries\footnotesize, inner sep=0pt] at (0,7.5) {440\,mL};
    \fill[bkRed] (-8,14) rectangle (8,23);
    \node[font=\sffamily\bfseries\footnotesize, text=white, inner sep=0pt] at (0,17.2) {10\%\R{増量}{ぞう|りょう}};
  \end{scope}
\end{tikzpicture}
\end{minipage}

\bkproblem （1）\R{去年}{きょ|ねん}480\R{人}{にん}だったお\R{祭}{まつ}りの\R{参加者}{さん|か|しゃ}が、\R{今年}{ことし}は552\R{人}{にん}でした。\R{今年}{ことし}は\R{去年}{きょ|ねん}の\R{何}{なん}\%ですか。（2）\qty{60}{W}の\R{電気}{でん|き}を\R{入}{い}れると、\qty{45}{W}\R{分}{ぶん}の\R{仕事}{し|ごと}をするモーターの\R{効率}{こう|りつ}は\R{何}{なん}\%ですか。
\bkpans{（1）\nobreak 115\%\quad（2）\nobreak 75\%}

\bkproblemgroup{チャレンジ}{2}

\bkproblem A\R{店}{てん}では3000\R{円}{えん}のゲームが25\%\R{引}{び}き、B\R{店}{てん}では\R{同}{おな}じゲームが2600\R{円}{えん}で、さらにその10\%\R{引}{び}きです。どちらの\R{店}{みせ}が\R{何}{なん}\R{円}{えん}\R{安}{やす}いですか。
\bkpans{A\R{店}{てん}2250\R{円}{えん}、B\R{店}{てん}2340\R{円}{えん}で、A\R{店}{てん}が90\R{円}{えん}\R{安}{やす}い}

\noindent\begin{minipage}[c]{0.6\linewidth}
\bkproblem \R{自転車}{じ|てん|しゃ}の\R{前}{まえ}のギアは\R{歯}{は}の\R{数}{かず}が48、\R{後}{うし}ろのギアは\R{歯}{は}の\R{数}{かず}が16です。タイヤは1\R{回転}{かい|てん}で\qty{2}{m}\R{進}{すす}みます。（1）ペダルを1\R{回転}{かい|てん}させると、タイヤは\R{何}{なん}\R{回転}{かい|てん}しますか。（2）ペダルを1\R{分}{ぷん}に60\R{回転}{かい|てん}させると、1\R{分}{ぷん}で\R{何}{なん}m\R{進}{すす}みますか。
\bkpans{（1）\nobreak 3\R{回転}{かい|てん}\quad（2）\nobreak $60\times3\times2=360$（m）}
\end{minipage}\hfill
\begin{minipage}[c]{0.37\linewidth}\centering
\begin{tikzpicture}[x=1mm, y=1mm]
  % 前は歯の数 48、後ろは歯の数 16（モジュール 0.4 mm）。チェーンは外側の接線
  \draw[bkInk!80, line width=1.4pt, line join=round]
    (2.048,9.379) -- (30.683,3.126) arc[start angle=77.68, end angle=-77.68, radius=3.2]
    -- (2.048,-9.379) arc[start angle=-77.68, end angle=-282.32, radius=9.6] -- cycle;
  \bkiifgear{(0,0)}{48}{0.4}{0}{bkGray!25}{bkInk!70}
  \bkiifgear{(30,0)}{16}{0.4}{0}{bkGray!25}{bkInk!70}
  \draw[bkInk!80, line width=1.4pt] (2.048,9.379) -- (30.683,3.126) (2.048,-9.379) -- (30.683,-3.126);
  \node[anchor=north, font=\sffamily\footnotesize, inner sep=0pt] at (0,-12) {\R{前}{まえ}　\R{歯}{は} 48};
  \node[anchor=north, font=\sffamily\footnotesize, inner sep=0pt] at (30,-6) {\R{後}{うし}ろ　\R{歯}{は} 16};
\end{tikzpicture}
\end{minipage}

\noindent\begin{minipage}[c]{0.6\linewidth}
\bkproblem \R{水平}{すい|へい}に\qty{500}{m}\R{進}{すす}むと\qty{35}{m}\R{上}{あ}がる\R{坂}{さか}のこう\R{配}{ばい}は\R{何}{なん}\%ですか。この\R{坂}{さか}を\R{水平}{すい|へい}に\qty{1200}{m}\R{進}{すす}むと、\R{何}{なん}m\R{上}{あ}がりますか。
\bkpans{$35\div500=0.07$ で7\%、$1200\times0.07=84$ で\qty{84}{m}}
\end{minipage}\hfill
\begin{minipage}[c]{0.37\linewidth}\centering
\begin{tikzpicture}[x=1mm, y=1mm]
  % こう配 7%（たても横と同じ縮尺）
  \fill[bkGreen!20] (0,0) -- (40,0) -- (40,2.8) -- cycle;
  \draw[bkInk, line width=1pt] (0,0) -- (40,2.8);
  \draw[bkInk!60, line width=0.6pt] (0,0) -- (40,0) -- (40,2.8);
  \node[anchor=north, font=\sffamily\footnotesize, inner sep=0pt] at (20,-1.8) {\R{水平}{すい|へい}に\qty{500}{m}};
  \node[anchor=west, font=\sffamily\footnotesize, inner sep=0pt] at (41.5,1.4) {\qty{35}{m}};
  \bkcar{(6,1.4)}{0.42}{bkBlue}
\end{tikzpicture}
\end{minipage}

\begin{furikaeri}
\begin{checklist}
\item 2つの\R{量}{りょう}の\R{関係}{かん|けい}を\R{比}{ひ}で\R{表}{あらわ}し、\R{等}{ひと}しい\R{比}{ひ}をつくれる。\hfill{\footnotesize\sffamily\color{bkTheme}→\pageref{bk:2-5:hi}ページ}
\item \R{比}{ひ}をかけ\R{算}{ざん}でつないで、\R{歯車}{は|ぐるま}の\R{回転}{かい|てん}を\R{求}{もと}められる。\hfill{\footnotesize\sffamily\color{bkTheme}→\pageref{bk:2-5:tsunagu}ページ}
\item \R{割合}{わり|あい}を\R{百分率}{ひゃく|ぶん|りつ}で\R{表}{あらわ}し、「20\%\R{引}{び}き」の\R{量}{りょう}を\R{求}{もと}められる。\hfill{\footnotesize\sffamily\color{bkTheme}→\pageref{bk:2-5:biki}ページ}
\end{checklist}
\end{furikaeri}

\bkfukidashi{delta}{「\R{何}{なに}をもとにするか」で、\R{同}{おな}じ\R{数}{かず}でも\R{意味}{い|み}が\R{変}{か}わるんだ。\R{次}{つぎ}の2-6では、はかった\R{数}{かず}の「ずれ」を\R{割合}{わり|あい}で\R{表}{あらわ}すよ。}
